\input{../tex-inputs/latex-leseansicht-vorspann}

\section[Arthur Schnitzler an Gustav Schwarzkopf, 2{[}3{]}. 1. 1925]{L04528 Arthur Schnitzler an Gustav Schwarzkopf, 2{[}3{]}. 1. 1925}
\nopagebreak\mylabel{L04528v}
\rehead{ }\normalsize\beginnumbering\briefempfaengerindex{Schwarzkopf, Gustav@\textsc{Schwarzkopf, Gustav}!zzzSchnitzler, Arthur@\emph{von Arthur Schnitzler}!1925-01-232@{2{[}3{]}. 1. 1925}|(be}
\toendnotes[C]{\smallbreak\pagebreak[2]}
\correspDesc{Versand  durch Arthur Schnitzler am 2[3]. 1. 1925 in St. Moritz
\newline{}Übermittlung  am 23. 1. 1925 in St. Moritz
\newline{}Erhalt  durch Gustav Schwarzkopf im Zeitraum [24. 1. 1925
                  – 28. 1. 1925?] in Wien}\toendnotes[C]{\smallbreak}
\Standort{DLA, A:Schnitzler, HS.1985.1.1897.}
\physDesc{Bildpostkarte, 265 Zeichen
\newline{}Handschrift: Bleistift, lateinische Kurrent
\newline{}Versand: Stempel: »\nobreak{}\oindex{St. Moritz@\textbf{St. Moritz}|pwk}St\textcolor{gray}{.}
                                       Moritz-Dorf, 23. 1. 25, 21\nobreak{}«.  }\toendnotes[C]{\smallbreak}\pstart{}{\pb}Hrn\pend{}\pstart{}Gustav Schwarzkopf\pend{}\pstart{}Wien I\oindex{I., Innere Stadt@\textbf{I., Innere Stadt}|pw}\pend{}\pstart{}Tiefer Graben 17\oindex{Wien@\textbf{Wien}!I., Innere Stadt@\textbf{I., Innere Stadt}!Tiefer Graben 17@\textbf{Tiefer Graben 17}|pw}\pend{}{\bigskip}
\pstart
           \noindent{}\centering{}{\pb}\textcolor{gray}{\textbf{St. Moritz-Dorf\oindex{St. Moritz@\textbf{St. Moritz}|pw} im Winter}}\pend
           \vspace{1em}
\pstart
           \noindent{}{\pb}Herzliche Grüße! Die \label{K_L04528-1v}\edtext{7 Vorlesungen\eventindex{Kultur- und Kongresszentrum Liederhalle@\textbf{Kultur- und Kongresszentrum Liederhalle}!Lesung von Die dreifache Warnung, Das Schicksal des Freiherrn von Leisenbohg, Große Szene, 9.1.1925@Lesung von Die dreifache Warnung, Das Schicksal des Freiherrn von Leisenbohg, Große Szene, 9.1.1925|pwv}\eventindex{Städtische Schauspiele Baden-Baden@\textbf{Städtische Schauspiele Baden-Baden}!Lesung von Das Tagebuch der Redegonda, Die letzten Masken, Weihnachts-Einkäufe, 11.1.1925@Lesung von Das Tagebuch der Redegonda, Die letzten Masken, Weihnachts-Einkäufe, 11.1.1925|pwv}\eventindex{Paulussaal@\textbf{Paulussaal}!Lesung von Das Tagebuch der Redegonda, Die letzten Masken, Große Szene, 12.1.1925@Lesung von Das Tagebuch der Redegonda, Die letzten Masken, Große Szene, 12.1.1925|pwv}\eventindex{Stadtcasino Basel@\textbf{Stadtcasino Basel}!Lesung von Das Tagebuch der Redegonda, Die letzten Masken, Große Szene, 13.1.1925@Lesung von Das Tagebuch der Redegonda, Die letzten Masken, Große Szene, 13.1.1925|pwv}\eventindex{Casino Bern@\textbf{Casino Bern}!Lesung von Die dreifache Warnung, Die letzten Masken, Der grüne Kakadu, 15.1.1925@Lesung von Die dreifache Warnung, Die letzten Masken, Der grüne Kakadu, 15.1.1925|pwv}\eventindex{Tonhalle St. Gallen@\textbf{Tonhalle St. Gallen}!Lesung von Die dreifache Warnung, Die letzten Masken, Zum großen Wurstel, 16.1.1925@Lesung von Die dreifache Warnung, Die letzten Masken, Zum großen Wurstel, 16.1.1925|pwv}\eventindex{Kongresshaus Zürich@\textbf{Kongresshaus Zürich}!Lesung von Das Tagebuch der Redegonda, Die letzten Masken, Zum großen Wurstel, 19.1.1925@Lesung von Das Tagebuch der Redegonda, Die letzten Masken, Zum großen Wurstel, 19.1.1925|pwv}}{\lemma{\textnormal{\emph{7 Vorlesungen}}}\Cendnote{\textnormal{Lesung von Die dreifache Warnung, Das Schicksal des Freiherrn von Leisenbohg, Große Szene, 9.1.1925. In: A. S.: \emph{Kulturveranstaltungen}, \url{https://schnitzler-kultur.acdh.oeaw.ac.at/pmb183497.html}; Lesung von Das Tagebuch der Redegonda, Die letzten Masken, Weihnachts-Einkäufe, 11.1.1925. In: A. S.: \emph{Kulturveranstaltungen}, \url{https://schnitzler-kultur.acdh.oeaw.ac.at/pmb183503.html}; Lesung von Das Tagebuch der Redegonda, Die letzten Masken, Große Szene, 12.1.1925. In: A. S.: \emph{Kulturveranstaltungen}, \url{https://schnitzler-kultur.acdh.oeaw.ac.at/pmb183512.html}; Lesung von Das Tagebuch der Redegonda, Die letzten Masken, Große Szene, 13.1.1925. In: A. S.: \emph{Kulturveranstaltungen}, \url{https://schnitzler-kultur.acdh.oeaw.ac.at/pmb183521.html}; Lesung von Die dreifache Warnung, Die letzten Masken, Der grüne Kakadu, 15.1.1925. In: A. S.: \emph{Kulturveranstaltungen}, \url{https://schnitzler-kultur.acdh.oeaw.ac.at/pmb183530.html}; Lesung von Die dreifache Warnung, Die letzten Masken, Zum großen Wurstel, 16.1.1925. In: A. S.: \emph{Kulturveranstaltungen}, \url{https://schnitzler-kultur.acdh.oeaw.ac.at/pmb183543.html}; Lesung von Das Tagebuch der Redegonda, Die letzten Masken, Zum großen Wurstel, 19.1.1925. In: A. S.: \emph{Kulturveranstaltungen}, \url{https://schnitzler-kultur.acdh.oeaw.ac.at/pmb183557.html}.}}}\label{K_L04528-1} sind glücklich vorbei,
               und ich will hier die obligate \label{K_L04528-2v}\edtext{Erkältung}{\lemma{\textnormal{\emph{Erkältung}}}\Cendnote{\textnormal{Vgl. A. S.: \emph{Tagebuch}, 23. 1. 1925.}}}\label{K_L04528-2}
               auscuriren. Ich schrieb Ihnen \label{K_L04528-3v}\edtext{von Freiburg\oindex{Freiburg im Breisgau@\textbf{Freiburg im Breisgau}|pw}}{\lemma{\textnormal{\emph{von Freiburg}}}\Cendnote{\textnormal{XXXX Auszeichnungsfehler: Dokument L04356 nicht gefunden.}}}\label{K_L04528-3} aus. Wollen Sie mir ein
               Wort hieher (\uline{Privathotel}\oindex{Privat-Hotel Badrutt@\textbf{Privat-Hotel Badrutt}|pw}) adressiren, wies bei Ihnen geht?\pend
           
\pstart
           Ihr{\\[\baselineskip]}\spacefill\mbox{A.}\pend
           \leftskip=0em{}\selectlanguage{ngerman}\endnumbering\briefempfaengerindex{Schwarzkopf, Gustav@\textsc{Schwarzkopf, Gustav}!zzzSchnitzler, Arthur@\emph{von Arthur Schnitzler}!1925-01-232@{2{[}3{]}. 1. 1925}|)be}\mylabel{L04528h}
\begin{anhang}
\end{anhang}\newcommand{\dateiname}{L04528}\newcommand{\titel}{Arthur Schnitzler an Gustav Schwarzkopf, 2[3]. 1. 1925}\newcommand{\editorInnen}{Herausgegeben von Jahnke, SelmaMüller, Martin Anton}\input{../tex-inputs/latex-leseansicht-abspann}
